\subsection{Checklist: New Assets}
\label{app:checklist_new_assets}

\paragraph{Question.}
Does the paper introduce new assets, such as code, datasets, benchmark
systems, trained models, generated data, figures, tables, or other artifacts
that are intended for release or reuse?

\paragraph{Answer (Yes/No/NA).}
Yes.

\paragraph{Justification.}
The paper introduces and uses several paper-specific assets. The primary asset
is the source code for sparse Koopman autoencoders, including sparse LISTA-style
encoders, dense and structured Koopman transitions, training entry points,
evaluation utilities, benchmark adapters, result collectors, statistical
summaries, and figure/table generation scripts. These assets are represented in
the current repository by the core package, command-line tools, batch scripts,
tests, dependency metadata, and paper source.

The paper also introduces a controlled 15-system two-dimensional multibasin
benchmark. The manuscript describes these systems as procedurally generated
with Claude Opus 4.5 and documents their system keys, generator families, basin
counts, stored-step convention, and analytic vector-field definitions. The
benchmark includes basin labels, basin counts, and attractor metadata only as
evaluation annotations. In the intended training and deployment setting, the
models are not given basin labels, basin counts, or trajectory-to-basin
assignments.

The paper additionally uses a 10-system Dysts $dt{\times}30$ forecasting
benchmark. The Dysts equations and metadata come from the external
\texttt{dysts} package version resolved in the project environment, documented
as \texttt{dysts} 0.96 in the paper appendix. These public Dysts systems are not
new assets introduced by this paper, but the stored-step protocol, cache
construction settings, split conventions, evaluation scripts, and collected
paper summaries are paper-specific generated artifacts.

Current repository search also identifies generated paper-facing figures,
LaTeX tables, CSV/JSON metric sidecars, manifests, and support-family codebook
artifacts under paths such as \texttt{docs/figures/neurips\_paper\_2026} and
its \texttt{\_tables} subdirectory. The training code writes best and latest
model checkpoint files named \texttt{checkpoint.pt} and \texttt{last.pt}. In
this checkout, \texttt{ls runs} reported no entries, but generated
\texttt{results/} artifacts include some \texttt{.pt} stage-two local-map and
checkpoint files from routed diagnostics. These checkpoint-like files are
derived experiment artifacts and should be treated separately from the minimal
paper release unless the authors explicitly select them for archival release.

\paragraph{Documentation, release, and anonymization status.}
\begin{itemize}
  \item \textbf{Code release.} The releasable code asset should include the
  core package, benchmark definitions and adapters, training and evaluation
  tools, paper-facing collection and plotting scripts, tests, the README,
  \texttt{pyproject.toml}, and \texttt{uv.lock}. The current repository includes
  an MIT license file, but the checkout's license metadata contains an author
  name; an anonymous-review release must either use venue-approved anonymous
  licensing metadata or defer identity-revealing attribution until the public
  de-anonymized release.

  \item \textbf{Benchmark release.} The 15 procedural multibasin systems should
  be released as explicit code and documentation, not as an opaque generated
  dataset. The release should preserve the disclosure that Claude Opus 4.5 was
  used for procedural benchmark generation. Generated training, validation, and
  test trajectories can be regenerated from the released vector fields and
  stored-step protocol. Evaluation-only basin metadata should be clearly marked
  as unavailable to training-time support selection.

  \item \textbf{Dysts protocol release.} The Dysts benchmark should be released
  as adapter code, system lists, pinned dependency metadata, stored-step
  settings, cache split conventions, and evaluation commands. Because the
  underlying Dysts systems are external public assets, the release should cite
  Dysts and respect its license rather than redistributing it as a new dataset.

  \item \textbf{Generated figures and tables.} The paper-facing figure files,
  LaTeX tables, and compact CSV/JSON summaries used by the manuscript should be
  released or archived with enough metadata to trace each display to the
  corresponding benchmark, model row, seed aggregation rule, and statistical
  test. Exploratory figures, failed pilot outputs, and scratch manifests are not
  required assets unless they support a reported claim.

  \item \textbf{Checkpoints and large generated files.} Exhaustive training
  runs, private scratch directories, large Dysts cache tensors, and all
  exploratory checkpoints do not need to be released if the code and protocols
  regenerate them. If any model checkpoints or second-stage local-map files are
  released, each file should be paired with its model configuration, system,
  seed, training budget, checkpoint-selection rule, evaluation command, and
  checksum.

  \item \textbf{Anonymization.} The anonymous artifact should remove or
  neutralize author names, local usernames, absolute private paths, private
  scratch locations, SLURM job logs, queue records, and other cluster-specific
  provenance that is not necessary for reproduction. No human-subject data,
  personal data, user logs, images of people, medical records, or proprietary
  datasets are introduced by the paper assets.
\end{itemize}

\paragraph{Caveats.}
This checklist item records that new assets exist; it does not by itself verify
that a final public archive, DOI, or anonymous review URL has been prepared.
Before submission, the authors should make an explicit release decision for
model checkpoints and generated cache tensors, because the paper can be
reproduced from code and protocols but reviewers may still value selected
paper-facing checkpoints for auditability. The distinction between benchmark
metadata used for evaluation and information unavailable during training should
remain visible in the release documentation. If future versions add real-world
datasets, human-generated measurements, proprietary data, or safety-critical
deployment artifacts, this answer should be revisited.
